\section{Choosing native repairs through a relative path target}
\label{sec:136-chosen-native-paths}

\paragraph{Archive and exact scope.}
This concerns cote136, physical pp.~117,119,121, footers116,118,120,
scans117,119,121.  The result applies to independent native blocks
whose input is zero on the entire first skeleton.  Their coherent
parameter images may be chosen.  It does not extend an arbitrary
previous choice of those images or an arbitrary attachment under
an old algebra.  Within this scope it removes the extra original
native-weight hypothesis of ARCH-136-HIGH-NATIVE-RESIDUE.

We use the checked statements ARCH-136-DERIVED-TAIL-COMPARISON,
ARCH-136-RELATIVE-COHERENT-MOORE,
ARCH-136-COHERENT-PD-REALIZATION, ARCH-136-PATHS, and the supplied
AQ-136-RELATIVE-H1, with its full normalized-primitive construction.
All maps preserve \(S=\Gamma_{\mathbb Z}(T)\) and its
divided powers.  The same proof works over \(\mathbb Z_{(p)}\).

\paragraph{The independent source blocks.}
Fix finitely many admissible targets \(C_j\).  For an index
\(\alpha\), choose \(n_\alpha\geq2\), an odd form degree
\(q_\alpha\geq1\), a positive weight \(m_\alpha\), and an
integer \(p_\alpha\geq2\).  Put
\[
 X_\alpha=\Delta[n_\alpha],\quad
 G_\alpha=\operatorname{sk}_1X_\alpha,\quad
 D_\alpha=p_\alpha(q_\alpha+1),\quad
 v_\alpha=p_\alpha m_\alpha.
                                                               \tag{136.CN1}
\]
Suppose every target row \((C_j)_*^{m_\alpha,q_\alpha}\)
is horizontally acyclic, and prescribe
\[
      a_{j\alpha}\in C_j^{m_\alpha,q_\alpha}(X_\alpha,G_\alpha).
                                                               \tag{136.CN2}
\]
There is no assumed lower bound on \(v_\alpha\) relative
to the vertical admissibility cutoff in degree \(D_\alpha\).
Choose a free resolution \(P_{\alpha,\bullet}\to I_{M_\alpha}\)
with generator weights at least \(M_\alpha\), where
\[
 M_\alpha+v_\alpha\geq\max_j W_{C_j}(D_\alpha)
                                                               \tag{136.CN3}
\]
is the bound from ARCH-136-DERIVED-TAIL-COMPARISON.

Let \(E_\alpha\) be the actual universal native coherent
algebra on the relative simplex module
\(\mathbb Z[X_\alpha]/\mathbb Z[G_\alpha]\).
It has an odd unit disk \(a,b\), \(da=b\), on every
relative basis simplex, with \(z=\gamma_{p_\alpha}(b)\),
\(e=a\gamma_{p_\alpha-1}(b)\), and the complete coefficient
resolution parameters \(U_i,V_i\) of
ARCH-136-COHERENT-PD-REALIZATION, including its bottom relations.
Thus \(de=p_\alpha z\).  Faces and degeneracies take
relative basis simplices to basis simplices or zero; the indicated
native relations are therefore actual simplicial relations.
The universal property is precisely a native section together
with all its coherent parameter maps satisfying those relations.
Set
\[
                         E=\bigotimes_{\alpha,S}E_\alpha.
                                                               \tag{136.CN4}
\]
Any set of blocks is allowed; elements have finite tensor support.
This is the strict PD coproduct of the independent blocks over
the fixed coefficients.  It has its evident augmentation
\(\epsilon:E\to S\), and \(E_0=E_1=S\).

\paragraph{ARCH-136-CHOSEN-NATIVE-PATHS.}
There are choices of all coherent target images and actual
strict simplicial maps
\[
        H_j:E\longrightarrow P_{\Delta[1]}C_j,
        \qquad e_0H_j=\eta_j\epsilon,                  \tag{136.CN5}
\]
such that \(e_1H_j\) sends every native generator to its
prescribed input (136.CN2).  At every simplex the augmentation
ideal of \(E\) is sent to sections zero on
\(\operatorname{sk}_1\Delta[n]\times\Delta[1]\).
For one target-relative comparison bound through degree three,
every closed degree-two Moore-cone pair in that
ideal has an image under \(e_1H_j\) with a primitive relative
to the whole first skeleton.  The bound is independent of the
modulus, simplex dimension, number of blocks, and their native
primes and arities.

\begin{proof}
Work with one block and one target.  Write \(I=\Delta[1]\),
\(X=X_\alpha\), \(G=G_\alpha\), and
\[
 Z=X\times I,\qquad Y=(G\times I)\cup(X\times\{0\}).
                                                               \tag{136.CN6}
\]
On \((X\times\partial I)\cup(G\times I)\), prescribe
the native input to be zero at endpoint zero and on \(G\times I\),
and to be \(a_{j\alpha}\) at endpoint one.  These prescriptions
agree on overlaps.  Acyclicity of the original input row extends
this section to \(\widetilde a\in C_j^{m,q}(Z,Y)\).
Indeed, normalized relative chains are free and split degreewise;
extend their chain map cell by cell, lifting each prescribed
boundary cycle in the acyclic row.  This is the same arbitrary-mono
extension argument used in ARCH-136-PATHS.

The relative pair \((Z,Y)\) is contractible: contracting the
interval to zero preserves \(Y\) and retracts both spaces onto
\(X\times\{0\}\).  Apply the derived-tail comparison
to \(C_j(Z,Y)\) and \(P_{\alpha,\bullet}[v]\).
Its total cohomology vanishes through degree \(D\) by (136.CN3).
The same is true of the required Moore cone in degree \(D-1\),
since that degree uses only the adjacent cohomology groups in
degrees \(D-1,D\).  The native pair
\[
 \widetilde z=\gamma_p(d\widetilde a),\qquad
 \widetilde e=\widetilde a\gamma_{p-1}(d\widetilde a),
 \qquad d\widetilde e=p\widetilde z
\]
therefore has the full coherent primitive on that coefficient
resolution.  Equivalently apply the already checked
ARCH-136-RELATIVE-COHERENT-MOORE to this pair.  It gives all
\(\widetilde U_i,\widetilde V_i\), including the bottom
relations, with values zero on \(Y\).

Sections on \(X\times I\) are sections on \(X\) of
the cotensor \(P_I C_j\).  The universal property of the
actual native coherent presentation gives
\(E_\alpha\to P_I C_j\).  Its endpoint zero is the
augmentation, endpoint one has the required native input, and
its whole first-skeleton restriction is the coefficient map.
This construction uses complete simplicial sections; it does
not choose incompatible values on individual faces.  Tensor
the maps for the independent blocks using multiplication in
the commutative PD target.  The maps agree on \(S\), so
this gives (136.CN5), also for arbitrary index sets.

Let \(B_n=\ker(E_n\to S)\).  Since \(E_0=E_1=S\),
\(B_n\) restricts to zero on the first skeleton.  A closed
pair \(\xi\in P_N(B_n)^2\) maps by \(H_j\) into
\[
             P_N C_j(\Delta[n]\times I,
                        \operatorname{sk}_1\Delta[n]\times I)^2.
                                                               \tag{136.CN7}
\]
Its endpoint zero is zero.  The bounded relative coefficient
comparison, followed by the cone of multiplication by \(N\),
identifies the cohomology in degree two in a uniform high-weight
range with ordinary relative cohomology with coefficients
\(S^w/NS^w\).  Only cohomology through degree three is
used; a bound dominating the underlying admissibility bounds
through form degree four suffices by the checked comparison.
It is independent of \(N\), exactly as in the derived
compatibility theorem.  Both endpoint inclusions of pairs are
homotopy inverses to projection onto \((\Delta[n],\operatorname{sk}_1\Delta[n])\).
They consequently induce the same map on that cohomology.
The endpoint-one class of (136.CN7) is zero.  It is thus an
actual cone boundary in the relative section complex, which
is the required primitive and proves individual residue vanishing.
\end{proof}

\paragraph{ARCH-136-CHOSEN-PATH-LOW-REPAIR.}
The algebra \(E\) and the maps (136.CN5) admit the supplied
RH1 iteration, producing \(E^+\), an augmentation
\(\epsilon^+:E^+\to S\), and strict maps
\[
 H_j^+:E^+\longrightarrow P_I C_j,
             \qquad e_0H_j^+=\eta_j\epsilon^+.
                                                               \tag{136.CN8}
\]
One still has \(E^+_0=E^+_1=S\); the integral vertical
\(H^0=S\), \(H^1=0\) conditions hold in one uniform
tail, and rationalization is unchanged by RH1.
The residue conclusion of ARCH-136-CHOSEN-NATIVE-PATHS holds
for every closed degree-two cone pair in
\(\ker(E_n^+\to S)\), including nonlinear classes created
or retained during the low repair.

\begin{proof}
Every native unit disk is rationally acyclic onto its constants.
The checked triangular coordinate calculation for the full
coherent presentation is a rational free unit-disk extension.
Consequently each \(E_\alpha\), and their arbitrary tensor
coproduct, has rational vertical cohomology exactly \(S_{\mathbb Q}\)
in every simplex and weight.  The same finite-support disk
contractions prove this for arbitrary ranks.  All weights of
new generators are positive, so weight zero is \(\mathbb Z[0]\).
ARCH-136-PATHS makes every \(P_I C_j\) admissible.  These
facts verify the hypotheses of AQ-136-RELATIVE-H1 with the
path algebras as retained targets.  Choose its one low-degree
bound to cover both the path targets and the original targets.
That theorem gives \(E^+\) and the extended paths.

Here are the extra augmentation and endpoint checks, which are
not consequences of horizontal admissibility of \(S\).
At one RH1 step, its normalized module of new zero-forms is
\(V_r=Z^1(N_r E)\) for \(r\geq1\) in high weights.
Send all those positive-normalized-degree generators to zero
under the augmentation.  This respects the boundary to
\(V_0=S\): the rational primitive in \(N_1E_{\mathbb Q}\)
has augmentation zero, and the coefficient boundary used to
define \(V_1\to S\) is consequently rationally zero.
The injection \(S\to S_{\mathbb Q}\) makes it integrally
zero.  Other normalized
boundaries are automatic.  The identifying zero-form relations
also have augmentation zero.  Thus the augmentation extends
through the actual PD pushout.  In unchanged low weights use
the old augmentation.

For endpoint zero, a newly adjoined normalized zero-form maps
under \(e_0H_j\) to a primitive of the zero one-form,
since the old map at that endpoint is the augmentation.
Normalized positive-degree \(H^0(C_j)\) is zero in the
chosen tail, and there are no negative forms.  The primitive
is therefore literally zero.  In normalized degree zero it
is the coefficient element.  This proves the endpoint identity
at every step, including its compatibility with every simplicial
operator.  Since \(E_1=S\), there are no normalized one-cycles
in original simplicial degree one, and the RH1 construction
does not alter that degree.  Induction and filtered colimits
prove (136.CN8) and \(E^+_0=E^+_1=S\).

Apply the last paragraph of the preceding proof to \(H_j^+\)
and its augmentation ideal.  The relative comparison bound is
still the same target comparison bound through degree three.  No assertion
that RH1 preserves a prior integral \(H^2\) or \(H^3\)
is required for this conclusion: every closed cone pair in the
resulting ideal has its own image along the retained path.
\end{proof}

\paragraph{ARCH-136-HIGH-DEGREE-PRESCRIBED-PATHS.}
For an independent block in the preceding construction, suppose
\(D\geq6\).  Its entire strict target map, including every
coherent \(U_i,V_i\), may already be prescribed.  Provided
that map is zero on the first skeleton and satisfies (136.CN3),
it still admits a path (136.CN5) from the augmentation.
The same holds simultaneously for any set of such blocks, and
the paths persist through the RH1 construction just proved.
There is no additional original native-weight guard.

\begin{proof}
Retain \(Z=X\times I\), but now set
\[
             Y'=(G\times I)\cup(X\times\partial I).
                                                               \tag{136.CN9}
\]
Interpolate the native input as before.  On \(Y'\) prescribe
all coherent parameters to be zero at endpoint zero and on
\(G\times I\), and to equal the given parameters at endpoint
one.  They agree on overlaps.  Let
\(\mathcal A=\operatorname{Tot}\operatorname{Hom}_S
(P_\bullet[v],C(Z))\), and let \(\mathcal R\) be its
restriction kernel for \(Y'\).  In every degree needed here,
restriction to \(Y'\) is onto: all \(P_i\) are free, and
all of their generator weights lie in the acyclic target rows
covered by (136.CN3).  This includes arbitrary products of
coordinate lifting problems.

The native pair gives a closed
\(\omega=(f_e,f_z)\in P_p(\mathcal A)^{D-1}\).
The prescribed boundary parameters give a cone primitive
\(\beta=(-V,-U)\) on \(Y'\); with
\(d(x,y)=(\delta x-py,-\delta y)\), one has
\(d\beta=\omega|_{Y'}\).  Lift \(\beta\) termwise
to \(\widetilde\beta\).  The only obstruction to correcting
that lift is the class
\[
 [\omega-d\widetilde\beta]\in
             H^{D-1}\bigl(P_p(\mathcal R)\bigr).
                                                               \tag{136.CN10}
\]
The derived-tail comparison and
ARCH-136-TAIL-EXT-VANISHING identify this group, in the stated
range, with
\[
                    H^{D-1}(Z,Y';S^v/pS^v).
                                                               \tag{136.CN11}
\]
Indeed the coefficient Hom complex is canonically
quasi-isomorphic to \(S^v[0]\), and its Moore cone has
only that quotient in degree zero.  The relative quotient of
(136.CN9) is the suspension of \(X/G\); equivalently its
relative normalized chains are homotopy equivalent to the
onefold shift of those of \((X,G)\).  Since \(X\) is a
simplex and \(G\) is a connected graph, the pair \((X,G)\)
has relative cohomology only in degree two.  The pair \((Z,Y')\)
therefore has relative cohomology only in degree three.
For \(D\geq6\), (136.CN11) is therefore zero.

Correct \(\widetilde\beta\) by a relative cone primitive
of (136.CN10).  The resulting full primitive gives all the
coherent cylinder parameters with exactly the prescribed
endpoint values.  The top resolution components of the same
cone equation impose both bottom relations; no separate lift
of a nonfree bottom module has been presumed.  The universal
presentation then gives the strict path.  Tensoring independent
paths and applying the preceding RH1 argument finish the proof.

For comparison, when \(D=4\) the possible obstruction group
is exactly \(H^1(G;S^v/pS^v)\).  This proof does not assert
that its particular class vanishes.  The chosen-map theorem
avoids prescribing that endpoint data, while the high-native
factorization gives a different sufficient condition for its
low residue.  These are distinct statements.
\end{proof}

\paragraph{Adding the common first skeleton.}
Suppose \(L(F)\) is the checked free even-disk left Kan
envelope of the common first skeleton, with fixed strict maps
\(L(F)\to C_j\).  Tensor it with \(E^+\), use the
constant path on \(L(F)\), and multiply with (136.CN8).
This gives a path from the composite
\[
 L(F)\otimes_S E^+\longrightarrow L(F)\longrightarrow C_j
\]
to the chosen target map.  The source has first skeleton \(F\).
Every closed cone pair in the ideal generated by
\(\ker(E^+\to S)\) has zero graph residue, by exactly
the same endpoint argument.  This statement is about that
ideal; it does not assert that all maps on the entire algebra
are augmentation paths to \(S\).

\paragraph{Exact consequence and remaining gate.}
The independent graph-zero native blocks, of arbitrary primes
and arities, can be assigned coherent images so that all their
subsequent low Moore residues vanish, even after their RH1
repair.  The native weights need only satisfy their original
input-row guard; the large parameter tails are allowed to vary
with the native operation.  This is a choice theorem for those
maps, not a theorem about every preassigned coherent map.

It does not apply directly to a horizontal attachment whose
boundary is already a nonzero element of an old source, or to
a native input with a nonzero common first-skeleton restriction.
A pointwise subtraction of an edgewise lift does not establish
a simplicially natural subtraction.  Nor may arbitrary further
horizontal cells be added while assuming that the retained
augmentation path persists.  If an admissible source retracted
onto \(L(F)\), its horizontal rows would force the same
acyclicity on that retract; this precludes inferring a full
horizontal completion from the displayed retraction alone.
The general integral common-source problem remains open.

\paragraph{Verification status.}
The author has completed a whole saved-proof reread, including
the row extension, contractible relative cylinder pair, full
coherent presentation, simultaneous simplicial coproduct,
uniform all-modulus relative comparison, augmentation through
the actual RH1 pushout, uniqueness at endpoint zero, unchanged
first skeleton, and the precise tensor-ideal limitation.
The parent and both peer branches completed whole saved-proof
checks of the first two statements and reported PASS.  The parent
reached the interpolation construction concurrently and contributed
the scope discussion; that reading is recorded as a contributor
review, not an independent rediscovery.  The path and RH1 assertions
are proved at the restricted independent-block scope.
The newly added prescribed-endpoint argument has also received
the author's complete proof reading, including the cone sign,
termwise free lifts, the suspension degree, and both bottom
relations.  The parent and reviewer123 subsequently read the whole saved
addendum and reported PASS.
